%% ============================================================
%% sections/s2-findings.tex
%% H. R. 9510 Bill v3.0 - SEC. 2. FINDINGS; THE EVIDENCE TO LAW RECORD.
%% Common visuals in this section: Figure 1 (four-work evidence-to-law lineage),
%% Table 1 (the four works at a glance), Figure 2 (accelerated timeline). The
%% findings are presented as a condensed narrative plus the three visuals, so the
%% section uses prose and figures in place of a long numbered list.
%% ============================================================
\billsec{SEC. 2.}{FINDINGS; THE EVIDENCE TO LAW RECORD.}

\uscprov{1}{\textbf{(a) In General.} Congress finds that an automated process
must verify robot-patient interaction code before that code is generated or
executed in a Physical AI oncology clinical investigation, and that the record of
the VVUQ-01 through VVUQ-04 developments, visualized in this section, establishes
both the need for and the feasibility of that requirement \cite{kawchak2026national,
vvuq05-figures-bill}.}

\figslot{Figure 1.}{Four-Work Evidence-to-Law Lineage}{VVUQ-01 method and
pipeline, to VVUQ-02 hard engineering proof, to VVUQ-03 proposed bill, to VVUQ-04
precise amendment; evidence flows forward and each work cites the records of the
prior works.}

\uscprov{1}{\textbf{(b) The Embodiment Problem and the Sequencing Gap.} Physical
AI oncology investigations place autonomous and semi-autonomous robots in direct
physical contact with patients, and the behavior of a robot at the moment of
contact is determined by software that is increasingly machine-generated. In a
disembodied pipeline an error yields a wrong number a human can catch downstream;
in an embodied surgical system the same error class yields a wrong motion a human
cannot catch in time. No Federal law requires the verification of robot-patient
interaction code to clear before that code is generated or executed, because the
order of operations is not regulated \cite{kawchak2026vvuq01, kawchak2026national}.}

\figslot{Table 1.}{The Four Works at a Glance}{Each work, what it is, its release
range, its headline assurance result (51 of 51 tests; 172 of 172 tests and
composite 93.56; 15 sections and 21 tables; new section 360e-5 and 10 conforming
amendments), and its end-goal DOI.}

\uscprov{1}{\textbf{(c) Feasibility Is Demonstrated.} Automated verification of
robot-patient interaction code is feasible today. A pipeline held the
verification fraction at 1.0, accepting one candidate artifact and blocking five,
and a ten-gate assurance suite over control software for a surgical humanoid
cleared reproducibly from a fixed seed, each gate bound to a published external
standard \cite{kawchak2026vvuq01, kawchak2026vvuq02}.}

\uscprov{1}{\textbf{(d) The Closest Existing Analogues.} The predetermined change
control plan authority in section 515C of the Federal Food, Drug, and Cosmetic
Act (21 U.S.C. 360e-4) and the Food and Drug Administration final guidance on such
plans for artificial-intelligence-enabled device software functions are the
closest existing analogues to validating a change before deployment, but each
remains voluntary, submission-based, and neither automated nor required in
advance of generation. The Quality Management System Regulation (part 820 of
title 21, Code of Federal Regulations), effective February 2, 2026, and the
electronic-records requirements of part 11 of such title supply quality and
record-integrity scaffolding but do not impose a rule that verification precede
generation \cite{fda-pccp-2024, fr-qmsr-2024}.}

\uscprov{1}{\textbf{(e) A Recognized Standards Basis.} Recognized consensus
standards, including ASME V\&V 40-2018, the Food and Drug Administration guidance
assessing the credibility of computational modeling, NASA-STD-7009A, IEEE Std
7009-2024, IEC 80601-2-77, ISO 14971, and ISO/TS 15066, supply the credibility
framework, dispersion bounds, fail-safe requirements, and biomechanical contact
limits to which an automated gate may be bound \cite{asmevv40, fda-credibility-2023,
nasastd7009, isots15066}.}

\uscprov{1}{\textbf{(f) Human-Over-AI Review Is Already Widely Adopted.}
Human-over-AI review is already widely adopted in State law, including the
California Physicians Make Decisions Act (Senate Bill 1120), the Illinois Wellness
and Oversight for Psychological Resources Act (House Bill 1806), the Texas
Responsible Artificial Intelligence Governance Act (House Bill 149), and Maryland
House Bill 820, and the strongest Federal precedent for mandatory human-over-AI
oversight is the Medicare Advantage artificial-intelligence guardrail (section
422.101(c) of title 42, Code of Federal Regulations) \cite{ca-sb1120, il-hb1806,
tx-hb149, md-hb820, cms-4201-ma}.}

\figslot{Figure 2.}{Accelerated Timeline Versus Conventional Methods}{The
conventional months-to-years path for a method paper, verified robotics code, a
manuscript, and a federal bill, set beside the about eleven days of autonomous,
real-time work that produced VVUQ-01 through VVUQ-04.}

\uscprov{1}{\textbf{(g) Modernization, Nondiscrimination, and Classification.} As
clinical investigation moves toward real-time and decentralized designs governed
by harmonized good clinical practice, manual review cannot keep pace with the
volume of machine-generated control code, and an automated gate that clears
before generation serves that modernization. Federal nondiscrimination law,
including section 92.210 of title 45, Code of Federal Regulations, title VI of the
Civil Rights Act of 1964 (42 U.S.C. 2000d), and the Americans with Disabilities
Act of 1990 (42 U.S.C. 12101 et seq.), requires that a patient-facing algorithm
minimize the risk of bias and respect accessibility, and the artificial-
intelligence taxonomy maintained by the American Medical Association provides a
graduated spine for matching verification rigor to the degree of autonomy of a
device \cite{iche6r3, cfr-92-210, usc-titlevi, usc-ada, cpt-appendix-s}.}

\begin{draftbox}{Section 2 - findings and the three findings visuals}
\dstep{Figure 1 (four-work evidence-to-law lineage): render the lineage diagram verbatim from papers/VVUQ-05/update-bill/figures-bill/01RootREADME.md (Visual 3, under the heading "Building process from one work to the next"), reconciled against the same diagram in papers/README.md and the "Lineage" diagram in papers/VVUQ-04/final-bill/README.md. Set it inside the asciifig environment so the boxes, arrows, and DOIs keep their exact spacing on a white background; place it as a full-width text figure. Do not redraw it as a raster image.}
\dstep{Table 1 (the four works at a glance): build the table from papers/VVUQ-05/update-bill/figures-bill/01RootREADME.md (Visual 2) and papers/README.md ("The four works at a glance"). Use the xltabular Y and L columns at the body measure with columns Work, What it is, Releases, Headline assurance result, and End-goal DOI; keep every cell left aligned and ragged right (each width prepended with \rabs) so no big interword gaps open; reword long cells onto two lines rather than letting any cell overflow.}
\dstep{Figure 2 (accelerated timeline): render the conventional-versus-this side-by-side timeline verbatim from papers/VVUQ-05/update-bill/figures-bill/01RootREADME.md (Visual 5) and papers/README.md ("Accelerated timeline versus conventional methods") in the asciifig environment; confirm the per-work spans (VVUQ-01 about 5 days, VVUQ-02 about 2 days, VVUQ-03 about 1 day, VVUQ-04 about 2 days, about eleven days total).}
\dstep{Findings prose: carry the fourteen grounded findings of Bill v2.0 from papers/VVUQ-04/final-bill/main.tex (SEC. 2, paragraphs (1) through (14)) and condense them into the lettered narrative paragraphs above, preserving every cited authority but using prose and the three visuals in place of a long numbered list, consistent with the instruction to use fewer numbered and lettered lists. Keep each finding tied to a recorded authority and resolve every citation to a key in references.bib.}
\dstep{Grounding: re-confirm the technical findings against papers/VVUQ-04/instruct-bill/09-vvuq-standards-clinical-trial-oncology.md and papers/VVUQ-03/final-paper/sections/definitions.tex; the legal-context findings against papers/VVUQ-04/instruct-bill/01-federal-statutory-framework.md, 02-fda-ai-device-regulation.md, 04-cms-coverage-payment-ai.md, and 05-privacy-security-nondiscrimination.md; and the state human-in-the-loop practice against 06-state-medical-ai-laws.md. Keep the 119th Congress emerging bills and the executive actions out of the findings and confine them to Appendix D per 10-legal-crosswalk-and-bill-style.md (Section B).}
\dstep{Formatting: single hyphens only; the section symbol for every codified reference; even RaggedRight spacing; reword so no finding ends in a one- or two-word line and no single line is stranded onto the next page; keep the white background on both figures and the table.}
\end{draftbox}
